% !TeX root = ../main.tex
\section{Wprowadzenie}

\begin{frame}{Przykład: tekst i graf}
  \begin{columns}[T,onlytextwidth]
    \begin{column}{0.5\textwidth}
      \begin{itemize}
        \item<1-> Sieć modelujemy jako graf $G = (V, E)$.
        \item<2-> Proces $p_1$ inicjuje rozgłaszanie.
        \item<3-> Wiadomość dociera do sąsiadów w kolejnych rundach.
        \item<4-> Liczba rund $\leq$ średnica grafu $D$.
      \end{itemize}
    \end{column}

    \begin{column}{0.45\textwidth}
      \centering
      \begin{tikzpicture}[node distance=14mm]
        \node[proc] (p1) {$p_1$};
        \node[proc, above right=of p1] (p2) {$p_2$};
        \node[proc, below right=of p1] (p3) {$p_3$};
        \node[proc, below right=of p2] (p4) {$p_4$};
        \node[proc, right=of p4] (p5) {$p_5$};

        \draw[link] (p1) -- (p2);
        \draw[link] (p1) -- (p3);
        \draw[link] (p2) -- (p4);
        \draw[link] (p3) -- (p4);
        \draw[link] (p2) -- (p3);
        \draw[link] (p4) -- (p5);

        % step-by-step highlighting
        \node<2->[proc, active] at (p1) {$p_1$};
        \draw<3->[msg] (p1) -- (p2);
        \draw<3->[msg] (p1) -- (p3);
        \node<3->[proc, active] at (p2) {$p_2$};
        \node<3->[proc, active] at (p3) {$p_3$};
        \draw<4->[msg] (p2) -- (p4);
        \draw<4->[msg] (p4) -- (p5);
        \node<4->[proc, active] at (p4) {$p_4$};
        \node<4->[proc, active] at (p5) {$p_5$};
      \end{tikzpicture}
    \end{column}
  \end{columns}
\end{frame}
